\subsection{CHARACTERS OF FINITE DIMENSIONAL g-MODULES}

Let \((e^{\lambda})_{\lambda \in \mathfrak{h}^{*}}\) be the canonical basis of the ring \(\mathbf{Z}[\mathfrak{h}^*]\) . Give the space \(\mathbf{Z}^{\mathfrak{h}^*}\) of all maps from \(\mathfrak{h}^*\) to \(\mathbf{Z}\) the product topology of the discrete topologies on the factors. If \(\varphi \in \mathbf{Z}^{\mathfrak{h}^*}\) , the family \((\varphi (\nu)e^{\nu})_{\nu \in \mathfrak{h}^*}\) is summable, and

\[
\varphi = \sum_ {\nu \in \mathfrak {h} ^ {*}} \varphi (\nu) e ^ {\nu}.
\]

Let \(\mathbf{Z}\langle\mathrm{P}\rangle\) be the set of \(\varphi\in\mathbf{Z}^{\mathfrak{h}^{*}}\) whose support is contained in a finite union of sets of the form \(\nu-\mathrm{P}_{+}\) , where \(\nu\in\mathfrak{h}^{*}\) . Then \(\mathbf{Z}[\mathrm{P}]\subset\mathbf{Z}\langle\mathrm{P}\rangle\subset\mathbf{Z}^{\mathfrak{h}^{*}}\) . Define on \(\mathbf{Z}\langle\mathrm{P}\rangle\) a ring structure extending that of \(\mathbf{Z}[\mathrm{P}]\) by putting, for \(\varphi,\psi\in\mathbf{Z}\langle\mathrm{P}\rangle\) and \(\nu\in\mathfrak{h}^{*}\) ,

\[
(\varphi \psi) (\nu) = \sum_ {\mu \in \mathfrak {h} ^ {*}} \varphi (\mu) \psi (\nu - \mu)
\]

(the family \((\varphi(\mu)\psi(\nu - \mu))_{\mu \in \mathfrak{h}^*}\) has finite support, in view of the condition satisfied by the supports of \(\varphi\) and \(\psi\) ). If \(\varphi = \sum_{\nu} x_{\nu} e^{\nu}\) and \(\psi = \sum_{\nu} y_{\nu} e^{\nu}\) , then \(\varphi \psi = \sum_{\nu, \mu} x_{\nu} y_{\mu} e^{\nu + \mu}\) .

Let \(\nu \in \mathfrak{h}^*\) . A partition of \(\nu\) into positive roots is a family \((n_{\alpha})_{\alpha \in \mathrm{R}_{+}}\) , where the \(n_{\alpha}\) are integers \(\geq 0\) such that \(\nu = \sum_{\alpha \in \mathrm{R}_{+}} n_{\alpha} \alpha\) . We denote by \(\mathfrak{P}(\nu)\) the number of partitions of \(\nu\) into positive roots. We have

\[
\mathfrak {P} (\nu) > 0 \Longleftrightarrow \nu \in \mathrm{Q} _ {+}.
\]

In this paragraph, we denote by K the following element of \(Z\langle P\rangle\) :

\[
\mathrm{K} = \sum_ {\gamma \in \mathrm{Q} _ {+}} \mathfrak {P} (\gamma) e ^ {- \gamma}.
\]

Now recall (Chap. VI, §3, no. 3, Prop. 2) that

\[
d = \prod_ {\alpha \in \mathrm{R} _ {+}} (e ^ {\alpha / 2} - e ^ {- \alpha / 2}) = \sum_ {w \in \mathrm{W}} \varepsilon (w) e ^ {w \rho}
\]

is an anti-invariant element of \(\mathbf{Z}[\mathrm{P}]\) .

Lemma 1. In the ring \(\mathbf{Z}\langle \mathrm{P}\rangle\) , we have \(\mathrm{K}\) . \(\prod_{\alpha \in \mathrm{R}_{+}}(1 - e^{-\alpha}) = \mathrm{Ke}^{-\rho}d = 1\) . Indeed,

\[
\mathrm{K} = \prod_ {\alpha \in \mathrm{R} _ {+}} \left(e ^ {0} + e ^ {- \alpha} + e ^ {- 2 \alpha} + \dots\right)
\]

SO

\[
K e ^ {- \rho} d = \prod_ {\alpha \in \mathrm{R} _ {+}} \left(1 + e ^ {- \alpha} + e ^ {- 2 \alpha} + \dots\right) \prod_ {\alpha \in \mathrm{R} _ {+}} \left(1 - e ^ {- \alpha}\right) = 1.
\]

Lemma 2. Let \(\lambda \in \mathfrak{h}^*\) . The module \(Z(\lambda)\) (§6, no. 3) admits a character that is an element of \(\mathbf{Z}\langle P\rangle\) , and we have \(d.\) ch \(Z(\lambda) = e^{\lambda + \rho}\) .

Let \(\alpha_{1},\ldots ,\alpha_{q}\) be distinct elements of \(\mathbb{R}_+\) . The \(X_{-\alpha_1}^{n_1}X_{-\alpha_2}^{n_2}\dots X_{-\alpha_q}^{n_q}\otimes 1\) form a basis of \(\mathrm{Z}(\lambda)\) (§6, Prop. 6 (iii)). For \(h\in \mathfrak{h}\) , we have

\[
\begin{array}{l} h. (X _ {- \alpha_ {1}} ^ {n _ {1}} X _ {- \alpha_ {2}} ^ {n _ {2}} \ldots X _ {- \alpha_ {q}} ^ {n _ {q}} \otimes 1) \\ \qquad = [ h, X _ {- \alpha_ {1}} ^ {n _ {1}} \ldots X _ {- \alpha_ {q}} ^ {n _ {q}} ] \otimes 1 + (X _ {- \alpha_ {1}} ^ {n _ {1}} \ldots X _ {- \alpha_ {q}} ^ {n _ {q}}) \otimes h. 1 \\ \qquad = (\lambda - n _ {1} \alpha_ {1} - \dots - n _ {q} \alpha_ {q}) (h) (X _ {- \alpha_ {1}} ^ {n _ {1}} \ldots X _ {- \alpha_ {q}} ^ {n _ {q}} \otimes 1). \end{array}
\]

Thus, the dimension of \(\mathrm{Z}(\lambda)^{\lambda - \mu}\) is \(\mathfrak{P}(\mu)\) . This proves that \(\operatorname{ch} \mathrm{Z}(\lambda)\) is defined, is an element of \(\mathbf{Z}\langle \mathrm{P} \rangle\) , and that

\[
\operatorname{ch} Z (\lambda) = \sum_ {\mu} \mathfrak {P} (\mu) e ^ {\lambda - \mu} = \mathrm{Ke} ^ {\lambda}.
\]

It now suffices to apply Lemma 1.

Lemma 3. Let M be a g-module which admits a character ch(M) whose support is contained in a finite union of the sets \(\mu-P_{+}\) . Let U be the enveloping algebra of g, Z the centre of U, \(\lambda_{0}\in h^{*}\) , and \(\chi_{\lambda_{0}}\) the corresponding homomorphism from Z to k ( \(\S8\) , Cor. 1 of Th. 2). Assume that, for all \(z\in Z\) , \(z_{M}\) is the homothety with ratio \(\chi_{\lambda_{0}}(z)\) . Let \(D_{M}\) be the set of \(\lambda\in W(\lambda_{0}+\rho)-\rho\) such that \(\lambda+Q_{+}\) meets Supp(ch M). Then ch(M) is a Z-linear combination of the ch \(Z(\lambda)\) for \(\lambda\in D_{M}\) .

If \(\mathrm{Supp(chM)}\) is empty, the lemma is clear. Assume that \(\mathrm{Supp(chM)} \neq \emptyset\) . Let \(\lambda\) be a maximal element of this support, and put \(\dim M^{\lambda} = m\) . There exists a \(\mathfrak{g}\) -homomorphism \(\varphi\) from \((Z(\lambda))^{m}\) to M which maps \((Z(\lambda)^{\lambda})^{m}\) bijectively onto \(M^{\lambda}\) ( \(\S 6\) , no. 3, Prop. 6 (i)). Thus, the central character of \(Z(\lambda)\) is \(\chi_{\lambda_0}\) , so \(\lambda \in W(\lambda_0 + \rho) - \rho\) ( \(\S 8\) , no. 5, Cor. 1 of Th. 2). This proves that \(D_M \neq \emptyset\) , and allows us to argue by induction on \(CardD_M\) . Let L and N be the kernel and cokernel of \(\varphi\) . Then we have an exact sequence of \(\mathfrak{g}\) -homomorphisms:

\[
0 \to \mathrm{L} \to (\mathrm{Z} (\lambda)) ^ {m} \to \mathrm{M} \to \mathrm{N} \to 0
\]

SO

\[
\operatorname{ch} (\mathrm{M}) = - \operatorname{ch} (\mathrm{L}) + m \operatorname{ch} \mathrm{Z} (\lambda) + \operatorname{ch} (\mathrm{N})
\]

(§7, no. 7, formula (6)). The sets Supp(ch L) and Supp(ch N) are contained in a finite union of sets \(\mu -\mathrm{P}_{+}\) . For \(z\in \mathbf{Z}\) , \(z_{\mathrm{L}}\) and \(z_{\mathrm{N}}\) are homotheties with ratio \(\chi_{\lambda_{0}}(z)\) . Clearly, \(D_{N} \subset D_{M}\) . On the other hand, \((\lambda + Q_{+}) \cap \text{Supp(ch M)} = \{\lambda\}\) , and \(\lambda \notin \text{Supp(ch N)}\) , so \(\lambda \notin D_{N}\) and

\[
\mathrm{Card} \mathrm{D} _ {\mathrm{N}} <   \mathrm{Card} \mathrm{D} _ {\mathrm{M}}.
\]

On the other hand, L is a submodule of \((Z(\lambda))^{m}\) ; if \(\lambda^{\prime} \in D_{L}\) , then \(\lambda^{\prime} + Q_{+}\) meets Supp(ch L) ⊂ Supp ch Z(λ), so \(\lambda \in \lambda^{\prime} + Q_{+}\) (§6, no. 1, Prop. 1 (ii)); it follows that \(D_{L} \subset D_{M}\) . Since \(L \cap (Z(\lambda)^{\lambda})^{m} = 0\) , we have \(\lambda \notin D_{L}\) , so

\(\mathrm{CardD_L} <   \mathrm{CardD_M}\)

It now suffices to apply the induction hypothesis.

THEOREM 1 (Character Formula of H. Weyl). Let M be a finite dimensional simple g-module, and \(\lambda\) its highest weight. Then

\[
\left(\sum_ {w \in \mathrm{W}} \varepsilon (w) e ^ {w \rho}\right). \mathrm{chM} = \sum_ {w \in \mathrm{W}} \varepsilon (w) e ^ {w (\lambda + \rho)}.
\]

With the notations of Lemma 3, the central character of M is \(\chi_{\lambda}\) ( \(\S6\) , no. 4, Prop. 7). Hence, by Lemmas 2 and 3, d.ch M is a Z-linear combination of the \(e^{\mu+\rho}\) such that

\[
\mu + \rho \in W (\lambda + \rho).
\]

On the other hand, by §7, no. 7, Lemma 7, d.ch M is anti-invariant, and its unique maximal term is \(e^{\lambda+\rho}\) , hence the theorem.

Example. Take \(\mathfrak{g} = \mathfrak{sl}(2,k)\) , \(\mathfrak{h} = kH\) . Let \(\alpha\) be the root of \((\mathfrak{g},\mathfrak{h})\) such that \(\alpha (H) = 2\) . The \(\mathfrak{g}\) -module \(V(m)\) has highest weight \((m / 2)\alpha\) . Hence

\[
\begin{array}{r l} & {\mathrm{ch} (\mathrm{V} (m)) = (e ^ {(m / 2) \alpha + \frac {1}{2} \alpha} - e ^ {- (m / 2) \alpha - \frac {1}{2} \alpha}) / (e ^ {\frac {1}{2} \alpha} - e ^ {- \frac {1}{2} \alpha})} \\ & {\qquad = e ^ {- (m / 2) \alpha}. (e ^ {(m + 1) \alpha} - 1) / (e ^ {\alpha} - 1)} \\ & {\qquad = e ^ {- (m / 2) \alpha} (e ^ {m \alpha} + e ^ {(m - 1) \alpha} + \dots + 1)} \\ & {\qquad = e ^ {(m / 2) \alpha} + e ^ {(m - 2) \alpha / 2} + \dots + e ^ {- (m / 2) \alpha}} \end{array}
\]

which also follows easily from §1, no. 2, Prop. 2.

